% ================================================================
%  conteudo/material_metodos.tex
% ================================================================

\section{Material e Métodos}

\subsection{Microrganismo}

O microrganismo utilizado foi a levedura \textit{Saccharomyces cerevisiae} ATCC~7754.

\subsection{Preparo do Meio de Cultura}

O cultivo foi conduzido em meio YM, cuja formulação consistiu em glicose (10~g/L), extrato de levedura (3~g/L), extrato de malte (3~g/L) e peptona (5~g/L). Após a homogeneização dos componentes, o pH da solução foi ajustado para 5,0.

\subsection{Cultivo}

A suspensão celular foi obtida a partir da inoculação da levedura \textit{Saccharomyces cerevisiae} em 150~mL do meio YM, seguida de incubação em agitador \textit{shaker} a 30~°C e 100~rpm, durante 27 horas.

\subsection{Determinação da Biomassa Seca}

Para quantificar a biomassa seca, utilizaram-se alíquotas de 10~mL da suspensão celular, acondicionadas em tubos de centrífuga previamente tarados. O material foi centrifugado a 3500~rpm durante 15~minutos e o sobrenadante foi descartado. Realizaram-se apenas uma lavagem com água destilada, seguidas de centrifugação para remoção de resíduos do meio de cultivo. Após o descarte do sobrenadante, os tubos contendo a biomassa foram levados à estufa a 105~°C até atingirem massa constante. A massa seca foi determinada conforme a Equação~\ref{eq:massa_seca} e a concentração de biomassa conforme a Equação~\ref{eq:concentracao}.

\subsection{Determinação da Absorbância}

A determinação da absorbância foi realizada a partir de diluições seriadas da suspensão celular, empregando-se os fatores de diluição de 5, 10, 15, 25, 50 e 100. As amostras foram preparadas em duplicata e submetidas à leitura de transmitância em espectrofotômetro a 600~nm. Os valores de transmitância (\%T) foram convertidos em absorbância pela Equação~\ref{eq:absorbancia}.

\subsection{Curva Padrão de Biomassa}

A curva padrão foi elaborada a partir da correlação entre a absorbância (turbidimetria a 600~nm) e as concentrações de biomassa determinadas pelo método da massa seca (g/L). Os dados foram utilizados para gerar um gráfico de concentração em função da absorbância, com ajuste linear e cálculo do coeficiente de determinação (R²).
